\documentclass{jarticle}
\usepackage{abstract}
\usepackage[dvipdfmx]{graphicx}

% ヘッダーを出力しない
\pagestyle{empty}

% ヘッダーを出力する場合は以下を設定
%\pagestyle{myheading}
%\kind{最終}% 中間 or 最終
%\year{4}% 年度
%\num{10}% No.

\title{電子工学科卒業研究発表会の予稿の書き方} %タイトル
\author{電子 太郎}% 報告者
\adviser{高専 太郎}% 指導教員

\begin{document}%

\maketitle % タイトルページの出力

\section{はじめに}%
どのような目的で本研究を実施するに至ったかを研究背景も
踏まえて説明（研究背景、解決すべき課題、研究目的など）する。
中間発表では、この1年で何をどこまで実施するかの研究計画について述べるとともに、
現時点での進捗状況、今後の課題について報告する。最終発表では、
中間発表以降に実施した研究内容を説明するとともに、
どのような研究成果が得られたかを具体的に説明すること。
  
\begin{figure}[b]
  \begin{center}
    \includegraphics[scale=1.0]{fig_1.bmp}
    \caption{未知抵抗の電圧電流特性測定回路}
  \end{center}
\end{figure}

\section{予稿の構成について}%
\subsection{研究方法}
理論的準備、実験装置の概略、実験方法などについて具体的に説明する。

\subsection{結果および考察}
実験結果とその考察を書く。実験項目が複数に分かれる場合は、
項目ごとに分けて書く。中間発表で研究結果がない場合は、
現在の進捗状況について説明する。

\subsection{まとめ}
実験の概略を説明して、その後に、
実験で得られた結果を数値を示して具体的に書く。
得られた結果が複数ある場合は、箇条書きにするなど、
わかり易く整理すること。中間発表の場合は、今後の課題についても書くこと。

\subsection{参考文献}
予稿の作成で論文、本、webページなどの
内容（文章、式、図、表など）を参考にし
た場合は参考文献に明記し、該当箇所に文献番号を付すこと。

\section{予稿の書き方の注意点}
\subsection{文章について}
主語と述語を明確に書くこと。接続詞副詞は、原則、
平がな書きとする。数字は、固有の専門用語、慣用句、
数詞を除いて、アラビア数字を用いること。
文中で変数を表す英字は、イタリック体を用いるのが
通例である\cite{ref:1}。ただし、
「単位」や「数学記号」はローマン体となるので注意すること。

\subsection{図表について}%
図表には、タイトルを付け（図は下、表は上）、
できる限り条件などの説明を加える。
本文を読まなくとも、図表だけを見て実験概略がわかるようにすること。
また、図表中の文字は小さすぎたり大きすぎたりしないように留意すること。

\begin{table}[t]
  \begin{center}
  \caption{未知抵抗の抵抗値の算出結果}
    \begin{tabular}{ccc} \Hline
    電圧$V$ [V]  & 電流$I$ [mA] & 抵抗値 [$\Omega$] \\ \hline
    1.0 & 3.0 &  333.3 \\
    2.0 & 6.1 &  327.9 \\  
    3.0 & 9.2 &  326.1 \\ \hline  
    \end{tabular}
  \end{center}
\end{table}

\begin{figure}[t]
  \begin{center}
    \includegraphics[scale=1.0]{fig_2.bmp}
    \caption{未知抵抗の電圧電流特性の測定結果}
  \end{center}
\end{figure}

\section{まとめ}
未知抵抗の電圧電流特性の電圧電流特性の測定を行った結果、
未知抵抗の抵抗値は、約330[$\Omega$]であることがわかった。

\begin{thebibliography}{99}%参考文献
\bibitem{ref:1}
著者 : ``タイトル'', 雑誌名,
 \textbf{号}, pp.開始ページ-終了ページ (発行年). 
\end{thebibliography}

\end{document}
